\cotefile{157-4}
\scanpage{243}
\mspage{117}
% archive sheet number: 242

\uline{Loc 4} (Axiome du carré cocartésien) Soit $X$ dans \unsure{$X$|$\mathcal{X}$}, $X_1$, $X_2$ deux ouverts de réunion $X$, $X_0=X_1\cap X_2$, d'où un diagramme cocartésien dans $\mathrm{Cat}$

\diagram{\begin{tikzcd}[ampersand replacement=\&,column sep=small,row sep=small]
 X_0 \arrow[r,"i_1"] \arrow[d,"i_2"'] \& X_1 \arrow[d,"j_2"] \\
 X_2 \arrow[r,"j_1"] \& \unsure{x|X}
\end{tikzcd}}

On a alors $i_1\in W\iff j_1\in W$ (et symétriquement $i_2\in W\iff j_2\in W$)

% slanted
\margin{Il y aura un Loc (5) \deleted{\unsure{plus}} ``basique'' plus bas, p.~124}

\uline{Loc} \textcircled{5} (Axiomes des $W$-fibrations) : \deleted{\unsure{la variante}}\inserted{il y a des} \deleted{\unsure{plus f}} variantes \textcircled{a} a$'$) \textcircled{b} b$'$) b$''$) c, c$'$, cf \uline{\unsure{XII}} p.~53-59. La forme la plus forte \textcircled{a} est celle-ci :

Soit $f:X\longrightarrow Y$ dans $\mathrm{Cat}$ \deleted{telle que pour toute}\inserted{une $W$-fibration} \deleted{\unsure{flèche} $Y''\to Y'$ dans} \unsure{au sens ancien} i.e. telle que pour toute $Y''\to Y'$ dans $\mathrm{Cat}_{/Y}$, avec $Y',Y''$ contractiles, \deleted{$X''=X\times_Y Y''$} la flèche $X''\to X'$ déduite par le chgt de base $X\to Y$ soit dans $W$. Alors $f$ est dans $\mathrm{Fib}_W$, i.e. le foncteur de chgt.\ de base \unsure{sur $f$} $\mathrm{Cat}_{/Y}\to\mathrm{Cat}_{/X}$ transforme flèches dans $W$ en flèches dans $W$.

\deleted{\uline{Loc \unsure{6|8}} (axiome de connexité) : \unsure{pour que} $\pi_0(f)$ \unsure{bijectif}}

% slanted
\margin{\deleted{Il y \unsure{avait} un Loc 6 concernant \unsure{les|des} \illegible{1} limites inductives}}

\uline{Loc} \deleted{\unsure{6|5}}7 (Axiome des limites inductives) $W$ est stable par limites inductives filtrantes

\inserted{forme renforcée de Loc 6}

\uline{Loc \deleted{\unsure{6|8}}7 bis}) (L'Accessibilité) $W$ est \deleted{accessible}\inserted{stable par $\mathrm{Lim}$} filtrantes \uline{et} accessible, i.e. il existe

\margin{\unsure{\deleted{la}|la|le}}

\uline{NB} \deleted{la conjonction de Loc 6 et Loc 7 \unsure{équivaut} à la \unsure{forme}} une petite partie \deleted{\unsure{des}} $W_{\unsure{\pi|0}}$ de $\mathrm{Fl}(\unsure{\mathrm{Cat}|\mathit{Cat}})$, telle que $W$ soit la plus petite partie de $\mathrm{Fl}(\mathrm{Cat})$ stable par \uline{lim} filtrantes, contenant $W_\pi$.
